% !TeX root = ../../Book4.tex
% 纲要标题：### 10.1 Yoneda 扩张
% 纲要标签：b4:sec:0901
% 纲要位置：计划纲要.md，第 3437 行。

\section{Yoneda 扩张}
\label{b4:sec:0901}

扩张不仅包含中间模，还包含核与商的指定映射。把扩张粘接起来得到的乘积，需要与消解中提升复形映射的乘积一致。

% 纲要内容：
%
% * Ext^1 与短正合扩张
% * 高阶扩张及等价
% * Yoneda 乘积与消解乘积的一致性
%

\subsection{\texorpdfstring{\(\operatorname{Ext}^1\) 与短正合扩张}{Ext1 与短正合扩张}}
\label{b4:subsec:0901:01}

给定左 \(R\)-模 \(M,N\)，\(\operatorname{Ext}_R^1(M,N)\) 既可由消解计算，也能描述把 \(N\) 放在核位置、把 \(M\) 放在商位置的所有短正合列。这里分类的数据包括两端的指定映射；仅知道中间模的同构类型还不够。

\begin{definition}\label{b4:def:yoneda-extension}
设 \(R\) 为含幺环，\(M,N\) 为幺性左 \(R\) 模。短正合列
\[
 \xi:\quad 0\longrightarrow N\xrightarrow{i}E\xrightarrow{p}M\longrightarrow0
\]
称为 \(M\) 被 \(N\) 的一个\term[kuozhang]{扩张}[extension]。给定另一个扩张 \(\xi':0\to N\xrightarrow{i'}E'\xrightarrow{p'}M\to0\)。若存在 \(R\) 模同态 \(f:E\to E'\)，满足 \(fi=i'\)、\(p'f=p\)，则称两扩张等价。这里两端的映射固定为恒等映射；两行短正合且两端为同构，由\cref{b4:thm:five-lemma}应用于包含两端零对象的五项序列，可知 \(f\) 自动为同构。仅有 \(E\cong E'\) 还不够，所用同构必须同时保持指定的核嵌入与商映射。
\end{definition}

\begin{theorem}\label{b4:thm:ext1-extensions}
设 \(R\) 为含幺环，\(M,N\) 为幺性左 \(R\) 模。\(M\) 被 \(N\) 的短正合扩张的等价类与 \(\operatorname{Ext}_R^1(M,N)\) 自然一一对应，分裂扩张对应零元。
\end{theorem}
\begin{proof}
选定正合列 \(0\to K\xrightarrow{j}P\xrightarrow{\varepsilon}M\to0\)，其中 \(P\) 投射。\cref{b4:thm:ext-long-exact}在第一变量给出的低次片段为
\[
 \operatorname{Hom}_R(P,N)\xrightarrow{j^*}\operatorname{Hom}_R(K,N)
 \longrightarrow\operatorname{Ext}_R^1(M,N)
 \longrightarrow\operatorname{Ext}_R^1(P,N)=0.
\]
最后一项由 \(P\) 的投射性消失，且 \(j^*(a)=aj\)，所以
\[
 \operatorname{Ext}_R^1(M,N)
 \cong\operatorname{Hom}_R(K,N)\big/\{\,a j:a\in\operatorname{Hom}_R(P,N)\,\}.
\]
分母恰好由可以从 \(K\) 延拓到 \(P\) 的映射组成。对扩张 \(\xi\)，由 \(P\) 投射且 \(p:E\to M\) 满射，可将 \(\varepsilon\) 提升为 \(t:P\to E\)，使 \(pt=\varepsilon\)。随后 \(ptj=\varepsilon j=0\)，故核映射 \(i:N\to E\) 的泛性质给出唯一的 \(u:K\to N\)，使 \(iu=tj\)。提升得到 \(t\)，经过核才得到 \(u\)；后者记录了 \(P\) 中的关系在扩张里落入 \(N\) 的方式。换成另一提升 \(t'\)，有 \(p(t'-t)=0\)，因而 \(t'-t=ia\)（\(a:P\to N\)），再由 \(i\) 单射得 \(u'-u=aj\)。扩张态射与 \(i,p\) 交换，将一套提升送到另一扩张的提升，故也保持这个余类。

反过来，希望把 \(P\) 中的 \(j(k)\) 与 \(N\) 中的 \(u(k)\) 识别，即令 \([(0,j(k))]=[(u(k),0)]\)。把两边相减便得到应商掉的关系 \((u(k),-j(k))\)，所以沿 \(u:K\to N\) 作推出得到
\[
 E_u=(N\oplus P)/\{\,(u(k),-j(k)):k\in K\,\}.
\]
由 \(n\mapsto[(n,0)]\) 及 \([(n,p)]\mapsto\varepsilon(p)\) 得到扩张。第一个映射单射，因为 \(j\) 单射；第二个映射的核由 \(\varepsilon(p)=0\) 即 \(p=j(k)\) 算出，此时 \([(n,j(k))]=[(n+u(k),0)]\)，故核恰好是 \(N\) 的像。如果 \(u-v=a j\)，考虑映射 \([(n,p)]\mapsto[(n+a(p),p)]\)。它把被商掉的关系送到
\[
 (u(k),-j(k))\longmapsto(u(k)-aj(k),-j(k))=(v(k),-j(k)),
\]
因而良定义；用 \(-a\) 得到逆映射，且两端映射保持不变。对原来的 \(\xi\)，映射 \([(n,p)]\mapsto i(n)+t(p)\) 也良定义，因为 \(iu=tj\)，并给出 \(E_u\to E\) 的扩张态射，因而是同构。两个构造互逆。

若类为零，则 \(u=aj\)（\(a:P\to N\)）。令 \(t'=t-ia\)，就有 \(pt'=\varepsilon\) 及
\[
 t'j=tj-iaj=iu-iu=0.
\]
因此 \(t'\) 唯一经过余核 \(\varepsilon:P\to M\)，写为 \(s\varepsilon\)；由 \(ps\varepsilon=\varepsilon\) 且 \(\varepsilon\) 满射，得到 \(ps=1_M\)。反之，若有截面 \(s\)，可取提升 \(t=s\varepsilon\)，它在 \(K\) 上为零，故对应零类。对 \(N\to N'\) 作推出相当于把 \(u\) 后复合，对 \(M'\to M\) 作拉回相当于提升该映射到投射消解后预复合，故对应自然。
\end{proof}

扩张类的加法也有直接构造。对两个扩张先取直和，再沿对角映射 \(M\to M\oplus M\) 拉回，使两个中间元素具有同一个商像；所得核仍有两份 \(N\)。随后沿加法 \(N\oplus N\to N\) 推出，把核元素 \((n,n')\) 合并为 \(n+n'\)。所得扩张称为\term[baerhe]{Baer 和}[Baer sum]。\footnote{拜尔（Reinhold Baer，1902-1979），德裔数学家，研究群的扩张、可解群和幂零群。}加法的核由 \((n,-n)\) 组成，故中间模具体为
\[
 (E\times_M E')/\{\,(i(n),-i'(n)):n\in N\,\}.
\]

\begin{proposition}\label{b4:prop:baer-addition}
设 \(R\) 为含幺环，\(M,N\) 为幺性左 \(R\) 模。Baer 和在扩张等价类上给出交换群结构；\cref{b4:thm:ext1-extensions}的对应是群同构。
\end{proposition}
\begin{proof}
从同一个 \(P\to M\) 分别提升到 \(E,E'\)，因两套提升的商像相同，可合并为到纤维积的提升；再作推出后，在 \(K\) 上诱导的映射就是 \(u+u'\)。因而 Baer 和对应 \(\operatorname{Ext}^1\) 中的加法。前一定理给出的双射于是把商群中已成立的结合律、交换律及零元、负元传到扩张类上，也证明结果不依赖代表元。具体地，分裂扩张为零，沿 \(-1_N\) 推出得到负元。
\end{proof}

例如，对素数 \(p\)，\(\operatorname{Ext}_{\mathbb Z}^1(\mathbb Z/p\mathbb Z,\mathbb Z/p\mathbb Z)=\mathbb Z/p\mathbb Z\)。非零类的中间群都同构于 \(\mathbb Z/p^2\mathbb Z\)，但指定的核嵌入可以不同，所以这些扩张仍有 \(p-1\) 个不同的非零类。这个区别在用扩张恢复结构时不能省略。

\subsection{高阶扩张及等价}
\label{b4:subsec:0901:02}

\begin{definition}\label{b4:def:higher-yoneda-extension}
设 \(R\) 为含幺环，\(M,N\) 为幺性左 \(R\) 模，\(n\ge1\) 为整数。给定左 \(R\) 模正合列
\[
 \xi:\quad 0\longrightarrow N\longrightarrow E_{n-1}\longrightarrow\cdots
 \longrightarrow E_0\longrightarrow M\longrightarrow0
\]
称为一个 \(n\) 阶扩张。两个这样的扩张之间，两端为恒等映射的链映射称为扩张态射。由扩张态射生成的等价关系称为 \term[yonedakuozhang]{Yoneda 扩张}[Yoneda extension]的等价关系：允许有限串扩张态射，箭头方向可以交替。这些扩张的等价类族暂记为 \(\operatorname{YExt}_R^n(M,N)\)；下面的定理将证明它可由一个集合参数化。
\end{definition}
\symidx{\(\operatorname{YExt}_R^n(M,N)\)：\(n\) 阶 Yoneda 扩张的等价类}

当 \(n>1\) 时，中间各项的映射不必可逆。例如令 \(A=\mathbb Z/2\mathbb Z\)，考察两个整数模二阶扩张
\[
\begin{aligned}
0&\longrightarrow A\xrightarrow{a\mapsto(a,0)}A\oplus\mathbb Z
 \xrightarrow{(a,z)\mapsto2z}\mathbb Z
 \xrightarrow{z\mapsto\bar z}A\longrightarrow0,\\
0&\longrightarrow A\xrightarrow{1_A}A\xrightarrow{0}A
 \xrightarrow{1_A}A\longrightarrow0.
\end{aligned}
\]
第一列中间的核为 \(A\oplus0\)，像为 \(2\mathbb Z\)，所以两列均正合。中间分别取投影 \((a,z)\mapsto a\) 和取模映射 \(z\mapsto\bar z\)，两端取恒等，就得到从第一列到第二列的扩张态射：中间方块两条复合均为零，其余方块也交换。中间项却分别是无限群与有限群，不可能同构。高阶等价因而不能要求逐项同构；要由这些不必可逆的态射得到等价关系，还须允许有限箭头串及箭头反向。下面同时说明这些类确实构成集合，并给出计算它们的方法。

\begin{theorem}\label{b4:thm:yoneda-ext-comparison}
设 \(R\) 为含幺环，\(M,N\) 为幺性左 \(R\) 模，\(n\ge1\) 为整数。扩张等价类可组成集合，并有对 \(M\) 反变、对 \(N\) 协变的自然双射
\[
 \operatorname{YExt}_R^n(M,N)\simeq\operatorname{Ext}_R^n(M,N).
\]
沿 \(M\) 的对角映射拉回、沿 \(N\) 的加法映射推出所定义的高阶 Baer 和，对应右侧的加法。
\end{theorem}
\begin{proof}
固定投射消解 \(P_\bullet\xrightarrow{\varepsilon}M\)，令 \(\Omega^nM=\ker(P_{n-1}\to P_{n-2})\)，其中 \(P_{-1}=M\)，\(P_0\to P_{-1}\) 为增广。记由微分诱导的满射为 \(\pi_n:P_n\to\Omega^nM\)。\cref{b4:thm:dimension-shift-basic,b4:prop:iterated-dimension-shift}的平移公式及短正合列 \(0\to\Omega^nM\to P_{n-1}\to\Omega^{n-1}M\to0\) 给出
\[
 \operatorname{Ext}_R^n(M,N)
 \simeq\operatorname{Hom}_R(\Omega^nM,N)
 /\operatorname{res}\operatorname{Hom}_R(P_{n-1},N).
\]
也可直接从消解读出此式。例如 \(n=2\) 时，\(\Omega^2M=\ker\mathrm{d}_{P,1}=\operatorname{im}\mathrm{d}_{P,2}\)。次数二的余圈 \(b:P_2\to N\) 满足 \(b\mathrm{d}_{P,3}=0\)，而 \(\ker\pi_2=\operatorname{im}\mathrm{d}_{P,3}\)，所以唯一写成 \(b=u\pi_2\)。次数二的余边是 \(a\mathrm{d}_{P,2}\)（\(a:P_1\to N\)），其下降映射恰为 \(a|_{\Omega^2M}\)。任意 \(n\ge1\) 同样有 \(\ker\pi_n=\operatorname{im}\mathrm{d}_{P,n+1}\)，因此余圈唯一经过 \(\Omega^nM\)，余边则对应商公式分母中的限制映射。这里沿用\secref{b4:sec:0802}投射消解计算中的无符号微分 \(f\mapsto f\mathrm{d}_P\)。

给定 \(n\) 阶扩张，将 \(E_0\to M\) 记为 \(\mathrm{d}_{E,0}\)，并将 \(P_0\to M\) 记为 \(\mathrm{d}_{P,0}\)。依次利用 \(P_i\) 的投射性，构造与微分、增广交换的提升 \(t_i:P_i\to E_i\)，\(0\le i<n\)。最后 \(t_{n-1}|_{\Omega^nM}\) 被下一微分零化，故经最左端的核嵌入 \(i:N\to E_{n-1}\) 分解，得到 \(u:\Omega^nM\to N\)。

为证明其余类与提升无关，只需把两套提升在最左端的差写成一个 \(P_{n-1}\to N\) 映射的限制。不能直接断言这一差落入 \(N\)：较低次数的提升也可能不同，须先逐层消去由它们传来的误差。对两套提升 \(t,s\)，从最低次数起构造 \(h_i:P_i\to E_{i+1}\)，使
\[
 t_i-s_i=\mathrm{d}_Eh_i+h_{i-1}\mathrm{d}_P\qquad(0\le i<n-1),\qquad h_{-1}=0.
\]
零次的差被增广零化。若已构造前一步，令
\[
 e_i=t_i-s_i-h_{i-1}\mathrm{d}_{P,i}.
\]
由两套提升的交换等式及上一同伦等式，得
\[
\begin{aligned}
\mathrm{d}_{E,i}e_i
&=(t_{i-1}-s_{i-1})\mathrm{d}_{P,i}
  -\mathrm{d}_{E,i}h_{i-1}\mathrm{d}_{P,i}\\
&=h_{i-2}\mathrm{d}_{P,i-1}\mathrm{d}_{P,i}=0
\qquad(i\ge1).
\end{aligned}
\]
所以尚未消去的误差落在核中；中间正合性把这个核识别为 \(E_{i+1}\) 的像，\(P_i\) 投射才使它提升为 \(h_i\)。到 \(i=n-1\) 时，核已经是 \(i(N)\)，于是存在 \(a:P_{n-1}\to N\) 使
\[
 t_{n-1}-s_{n-1}-h_{n-2}\mathrm{d}_{P,n-1}=ia.
\]
限制到 \(\Omega^nM\) 后，微分项为零，再由 \(i\) 单射得到 \(u_t-u_s=a|_{\Omega^nM}\)，这恰在商公式的分母中。\(n=1\) 时直接用 \(t_0-s_0\) 即可。扩张态射把一套提升变成另一扩张的提升，因而整个箭头串都保持该余类。

对任意 \(u:\Omega^nM\to N\)，把截短消解
\[
 0\longrightarrow\Omega^nM\longrightarrow P_{n-1}\longrightarrow\cdots
 \longrightarrow P_0\longrightarrow M\longrightarrow0
\]
的最左端沿 \(u\) 推出。新项是
\(E_{n-1}(u)=(N\oplus P_{n-1})/\{(u(x),-x)\}\)，其余项保持 \(P_i\)。与一阶情形相同，核及像的计算证明新列正合。若 \(u-v=a|_{\Omega^nM}\)，在新项上用 \([(z,p)]\mapsto[(z+a(p),p)]\)，在其他项上用恒等映射，就得到两个新扩张的同构。

从任意原扩张取出的 \(u\) 又给出这个推出扩张到原扩张的态射：最左新项用 \([(z,p)]\mapsto i(z)+t_{n-1}(p)\)，其他项用 \(t_i\)。标准推出扩张的这些项取自 \(P_\bullet\)，原扩张的 \(E_i\) 却不必投射，所构造的中间箭头也不必可逆；它们之所以足以识别两扩张，正是因为高阶等价由扩张态射生成。因此两个方向互逆。特别地，每个扩张类都有一个由 \(u\in\operatorname{Hom}_R(\Omega^nM,N)\) 得到的代表；这个 Hom 集的上述商集便给出全部扩张类的集合参数化。

接着核对两变量的作用。若 \(v:N\to N'\)，沿 \(v\) 推出扩张，并把原来的提升复合到推出列中，最左端诱导的映射变成 \(vu:\Omega^nM\to N'\)。它对应消解余圈的后复合，故与 \(\operatorname{Ext}_R^n(M,v)\) 相容。

若 \(f:M'\to M\)，取投射消解 \(P'_\bullet\to M'\)，由\cref{b4:thm:resolution-comparison}选取提升 \(f\) 的比较映射 \(c:P'_\bullet\to P_\bullet\)。记
\[
 c_\Omega=c_{n-1}|_{\Omega^nM'}:\Omega^nM'\longrightarrow\Omega^nM.
\]
这里 \(f\) 决定要拉回的扩张，\(c\) 实现 Ext 计算中的预复合，\(c_\Omega\) 则记录它在关系模上的作用。拉回列的最低项为 \(E_0\times_M M'\)，其提升明确取为
\[
 P'_0\longrightarrow E_0\times_M M',\qquad
 x\longmapsto(t_0c_0(x),\varepsilon_{P'}(x));
\]
它落在纤维积中，因为 \(\mathrm{d}_{E,0}t_0c_0=f\varepsilon_{P'}\)。其余次数用原提升与 \(c_i\) 的复合，最左端映射因而成为 \(u c_\Omega\)。若余圈 \(b:P_n\to N\) 写为 \(b=u\pi_n\)，记相应满射为 \(\pi'_n:P'_n\to\Omega^nM'\)，则
\[
 b c_n=u\pi_nc_n=u c_\Omega\pi'_n,
\]
所以这个商类恰是 \(\operatorname{Ext}_R^n(f,N)\) 所给的类。若换另一比较提升 \(c'\)，由\cref{b4:prop:comparison-homotopy-unique}，存在同伦 \(h_i:P'_i\to P_{i+1}\)，使
\[
 c_i-c'_i=\mathrm{d}_{P,i+1}h_i+h_{i-1}\mathrm{d}_{P',i},\qquad h_{-1}=0.
\]
在 \(\Omega^nM'\) 上第二项为零，因此
\[
 u c_\Omega-u c'_\Omega
 =\left.(u\pi_nh_{n-1})\right|_{\Omega^nM'}.
\]
其中 \(h_{n-1}:P'_{n-1}\to P_n\)，故 \(u\pi_nh_{n-1}\) 定义在整个 \(P'_{n-1}\) 上，右侧恰属于商公式的分母。这既证明拉回与 Ext 的预复合相容，也证明更换比较提升不改变所得类。

最后把换消解与换扩张提升区分开来。若 \(P'_\bullet\) 也是 \(M\) 的投射消解，上面的构造取 \(f=1_M\)，就把两种商描述按 \([u]\mapsto[u c_\Omega]\) 比较。取反向比较后，两种复合分别与相应的恒等映射同伦，所以诱导互逆的同构；余圈等式也表明这正是\chapref{b4:ch:08}用于识别两种 Ext 计算的比较同构。因而不同消解给出的识别在该同构下相容。前面固定消解时已经证明换扩张中的提升、沿扩张态射改变代表均不改变商类，这里消去的则是投射消解本身的选择。

对扩张 \(\xi,\xi'\)，先逐项取直和 \(\xi\oplus\xi'\)，得到右端为 \(M\oplus M\)、左端为 \(N\oplus N\) 的列；再在右端沿 \(\Delta_M:M\to M\oplus M\) 拉回，最后在左端沿 \(+_N:N\oplus N\to N\) 推出。也就是
\[
 \xi+\xi'=(+_N)_*\Delta_M^*(\xi\oplus\xi').
\]
分别取 \(P_\bullet\) 到两列的提升，再用逐项对角映射 \(P_\bullet\to P_\bullet\oplus P_\bullet\) 提升到拉回列。在最左端，所得映射是 \((u,u'):\Omega^nM\to N\oplus N\)；推出后成为 \(+_N(u,u')=u+u'\)。所以 Baer 和对应商集中的加法。特别地，它不依赖代表扩张，并由此给出与 Ext 相同的交换群结构。
\end{proof}

对 \(n=0\)，约定 \(\operatorname{YExt}_R^0(M,N)=\operatorname{Hom}_R(M,N)\)。在没有足够投射或内射对象的一般交换范畴中，仍可讨论扩张及其箭头串；本节已证的与导出函子的比较则是在模范畴中进行的。

\subsection{Yoneda 乘积与消解乘积}
\label{b4:subsec:0901:03}

若一列以 \(N\) 为商，另一列以 \(N\) 为核，就能把它们首尾连接。例如两个一次扩张
\[
0\longrightarrow L\xrightarrow{i}E\xrightarrow{v}N\longrightarrow0,
\qquad
0\longrightarrow N\xrightarrow{j}F\xrightarrow{p}M\longrightarrow0
\]
拼接为
\[
0\longrightarrow L\xrightarrow{i}E\xrightarrow{jv}F
\xrightarrow{p}M\longrightarrow0.
\]
由于 \(j\) 单射，\(\ker(jv)=\ker v=\operatorname{im}i\)；由于 \(v\) 满射，\(\operatorname{im}(jv)=j(N)=\ker p\)，所以拼接列仍正合。中间箭头是明确的复合 \(E\to N\to F\)，而不是将两列中的 \(N\) 直接略去。这一操作给出 Ext 群之间并非来自普通数乘的乘法。

\begin{definition}\label{b4:def:yoneda-product}
设 \(R\) 为含幺环，\(L,N,M\) 为幺性左 \(R\) 模，\(m,n\ge1\) 为整数。对 \(\alpha\in\operatorname{YExt}_R^m(N,L)\)、\(\beta\in\operatorname{YExt}_R^n(M,N)\)，把代表扩张在共同的 \(N\) 处拼接，删除两处 \(N\)，以原来的满射到 \(N\) 与从 \(N\) 出发的单射之复合作为中间箭头。所得 \((m+n)\) 阶扩张类记为 \(\alpha\beta\)，称为\term[yonedachengji]{Yoneda 乘积}[Yoneda product]。约定 \(\operatorname{YExt}_R^0(X,Y)=\operatorname{Hom}_R(X,Y)\)：左因子为零次同态 \(N\to L\) 时沿该同态推出，右因子为零次同态 \(M\to N\) 时沿该同态拉回；两个因子均为零次时取同态复合。
\end{definition}
\symidx{\(\alpha\beta\)：Yoneda 扩张类的乘积}

\begin{proposition}\label{b4:prop:yoneda-resolution-product}
设 \(R\) 为含幺环，\(L,N,M\) 为幺性左 \(R\) 模。Yoneda 乘积良定义、双加性且结合，并在 Yoneda 扩张与 Ext 的自然对应下等于下述消解乘积。取投射消解 \(P_\bullet\to M\)、\(Q_\bullet\to N\)，以 \(\varepsilon_Q:Q_0\to N\) 表示增广，设 \(m,n\ge0\) 为整数。若 \(b:P_n\to N\)、\(a:Q_m\to L\) 分别是 \(\operatorname{Hom}_R(P_\bullet,N)\)、\(\operatorname{Hom}_R(Q_\bullet,L)\) 中的余圈，把 \(b\) 逐次提升为
\[
 t_i:P_{n+i}\longrightarrow Q_i,\qquad
 \varepsilon_Qt_0=b,\qquad \mathrm{d}_Qt_i=t_{i-1}\mathrm{d}_P\quad(i\ge1),
\]
则乘积由余圈 \(a t_m:P_{n+m}\to L\) 代表。
\end{proposition}
\begin{proof}
拼接在共同项处正合，因为像和核都由被删除的 \(N\) 描述。若一个因子沿某个扩张态射改变，把该态射与另一列的恒等映射拼接，仍得到扩张态射。因此拼接保持等价关系的每个生成关系，也就保持由正向或反向扩张态射组成的任意有限箭头串。零次因子的拉回或推出也把扩张态射送到扩张态射，故这些情形同样良定义。三个正次数因子依次拼接不改变项及其复合箭头。零次因子位于外端时，拉回或推出的泛性质使两种顺序相容；它位于中间时，拉回的投影与推出的结构映射给出两种拼接列之间的扩张态射。连续的零次因子则由同态复合的结合律处理。因而乘积结合。

对于消解公式，先沿满射 \(\varepsilon_Q:Q_0\to N\) 提升余圈 \(b:P_n\to N\)，得到 \(t_0:P_n\to Q_0\)。因此 \(Q\) 的零次已经对应 \(P\) 的第 \(n\) 次。由 \(b\mathrm{d}_{P,n+1}=0\)，\(t_0\mathrm{d}_{P,n+1}\) 落在 \(\ker\varepsilon_Q\)，可用 \(P_{n+1}\) 的投射性提升至 \(Q_1\)。以后每步同样先用 \(\mathrm{d}_P^2=0\) 核对落在核中，再提升；每向 \(Q\) 的高次推进一步，源也从 \(P_{n+i}\) 推进到 \(P_{n+i+1}\)。到第 \(m\) 步，恰可形成
\[
 P_{n+m}\xrightarrow{t_m}Q_m\xrightarrow{a}L,
 \qquad
 a t_m\mathrm{d}_{P,n+m+1}
 =a\mathrm{d}_{Q,m+1}t_{m+1}=0.
\]
这既说明所得是余圈，也说明乘积次数为何是 \(n+m\)。

对 \(m,n\ge1\)，要比较两种乘积，只须读取拼接扩张最左端的关系映射，证明其余圈为 \(at_m\)。用\cref{b4:thm:yoneda-ext-comparison}的推出构造代表 \(\alpha,\beta\)：把 \(\alpha\) 的最左新项记为 \(A\)，把 \(\beta\) 的最左新项记为 \(B\)，其余项来自两条消解。设 \(a,b\) 经合冲模的分解分别为
\[
 \bar a:\Omega^mN\to L,\qquad \bar b:\Omega^nM\to N,
\]
则这两个推出项为
\[
\begin{aligned}
 A&=(L\oplus Q_{m-1})/\{(\bar a(x),-x):x\in\Omega^mN\},\\
 B&=(N\oplus P_{n-1})/\{(\bar b(y),-y):y\in\Omega^nM\}.
\end{aligned}
\]
以 \(i_\alpha:L\to A\)、\(i_\beta:N\to B\) 表示单射，以 \(j_\alpha:Q_{m-1}\to A\)、\(j_\beta:P_{n-1}\to B\) 表示 \(q\mapsto[(0,q)]\) 型的映射。要将两段提升在共同的 \(N\) 处接起来，须使经 \(P_{n-1}\) 进入 \(B\) 与经 \(Q_0\to N\) 进入 \(B\) 的两条路径相同，所需等式正是下面的方块：
\[
\begin{tikzcd}[column sep=large,row sep=large]
 P_n \ar[r,"\mathrm{d}_{P,n}"] \ar[d,"t_0"']
   & P_{n-1} \ar[d,"j_\beta"] \\
 Q_0 \ar[r,"i_\beta\varepsilon_Q"'] & B.
\end{tikzcd}
\]
它交换，因为 \(j_\beta\mathrm{d}_{P,n}=i_\beta b=i_\beta\varepsilon_Qt_0\)。因此从 \(P_0\) 到 \(P_{n-2}\) 的恒等映射、\(P_{n-1}\to B\) 的 \(j_\beta\)，以及
\[
 P_{n+i}\xrightarrow{t_i}Q_i\quad(0\le i<m-1),\qquad
 P_{n+m-1}\xrightarrow{j_\alpha t_{m-1}}A
\]
组成到拼接列的逐项提升；当 \(m=1\) 或 \(n=1\) 时，相应的指标区间为空。若 \(m=1\)，拼接箭头 \(A\to B\) 由 \(i_\beta\varepsilon_Q:Q_0\to B\) 及 \(L\to B\) 的零映射在推出上诱导，所以上面的方块仍给出所需相容性。

这些映射给出到拼接列的提升以后，最左端的关系映射由 \(P_{n+m}\to P_{n+m-1}\to A\) 读出。下面的上方方块使用提升等式，下方方块使用 \(A\) 的推出关系，将这条路径与 \(P_{n+m}\to Q_m\to L\to A\) 比较：
\[
\begin{tikzcd}[column sep=large,row sep=large]
 P_{n+m} \ar[r,"\mathrm{d}_{P,n+m}"] \ar[d,"t_m"']
   & P_{n+m-1} \ar[d,"t_{m-1}"] \\
 Q_m \ar[r,"\mathrm{d}_{Q,m}"] \ar[d,"a"']
   & Q_{m-1} \ar[d,"j_\alpha"] \\
 L \ar[r,"i_\alpha"'] & A.
\end{tikzcd}
\]
下方方块交换，是因为推出关系给出 \(j_\alpha\mathrm{d}_{Q,m}=i_\alpha a\)。于是
\[
 j_\alpha t_{m-1}\mathrm{d}_{P,n+m}=i_\alpha a t_m.
\]
由于 \(\operatorname{im}\mathrm{d}_{P,n+m}=\Omega^{n+m}M\) 且 \(i_\alpha\) 单射，拼接列在最左端诱导的映射正是余圈 \(a t_m\) 经 \(\Omega^{n+m}M\) 的分解。因此拼接扩张对应所陈述的消解乘积。该对应也证明换 \(t_i\)、换余圈代表或换消解均不改变所得 Ext 类。

固定 \(b\) 的提升时，\((a+a')t_m=a t_m+a't_m\)；固定 \(a\) 时，两套 \(b,b'\) 的提升之和提升 \(b+b'\)。因此乘积双加性。若 \(m=0\)，写 \(a=h\varepsilon_Q\)，其中 \(h:N\to L\)，则 \(a t_0=h b\)，正是沿 \(h\) 推出的后复合。若 \(n=0\)，写 \(b=f\varepsilon_P\)，其中 \(f:M\to N\)，则 \(t_i\) 是提升 \(f\) 的比较映射，\(a t_m\) 正是自然性中沿 \(f\) 拉回的余圈。两个次数均为零时，这又化为同态复合。
\end{proof}

于是 \(\bigoplus_{n\ge0}\operatorname{Ext}_R^n(M,M)\) 成为带单位的分次环，单位是 \(1_M\)。这里没有一般的分次交换性断言：次数零部分已经可能是非交换环 \(\operatorname{End}_R(M)\)。例如在域 \(R=k\) 上取 \(M=k^2\)，这一部分就是矩阵环 \(M_2(k)\)。

\ProseParagraphBreak
即使次数零部分交换，正次数也未必满足分次交换律。取域 \(k\) 及交换环 \(R=k[\epsilon]/(\epsilon^2)\)，令 \(M=k=R/(\epsilon)\)。\secref{b4:sec:0701}已证明它有各项为 \(R\)、微分均为乘 \(\epsilon\) 的周期自由消解；\secref{b4:sec:0802}由施加 Hom 后微分为零得到每次 Ext 为 \(k\)。次数一类 \(t\) 由自然商映射 \(q:P_1=R\to k\) 表示，提升可取底层均为 \(1_R\) 的 \(t_i:P_{i+1}\to P_i\)。它们使次数降低一，而不是原消解中的零次恒等映射。上述乘积公式将 \(t^2\) 表示为 \(q t_1:P_2\to k\)，它仍是非零商映射；Hom 复形的微分为零，故它也不是余边。因此 \(t^2\ne0\)。若 \(\operatorname{char}k\ne2\)，分次交换律却会要求 \(t^2=-t^2\)，这说明正次数中的障碍也确实存在。\cref{b4:ex:09-dual-number-product}进一步计算这一例子的整个自扩张环。

\ProseParagraphBreak
消解乘积的计算采用 \(\delta f=f\mathrm{d}\) 的次数约定。与\cref{b4:def:hom-complex}的内部 Hom 约定比较时，把 \(P_i,Q_i\) 放在上链次数 \(-i\)，并记 \(c_r=(-1)^{r(r+1)/2}\)。\secref{b4:sec:0802}的重标度将内部 Hom 的余圈 \(c_nb\) 对应于这里的 \(b\)。相应的次数 \(n\) 提升在分量 \(P_{n+i}\to Q_i\) 上须取
\[
 T_i=c_n(-1)^{ni}t_i,\qquad
 \mathrm{d}_QT_i=(-1)^n\cdot T_{i-1}\mathrm{d}_P\quad(i\ge1),
\]
才能满足内部 Hom 的余圈条件。将另一个余圈写成 \(c_ma\)，其复合为 \(c_mc_n(-1)^{mn}at_m\)。由于
\[
 c_{m+n}=(-1)^{mn}c_mc_n,
\]
在次数 \(m+n\) 作回同一重标度后，仍得到 \(at_m\)。两种符号约定因此给出同一个乘积，但须同时调整余圈与全部提升分量。Ext 扩张的基本对应和 Baer 和可参阅 Weibel\cite[Theorem 3.4.3, Definition 3.4.4, Corollary 3.4.5]{Weibel1994HomologicalAlgebra}；高阶扩张的进一步讨论见该书\cite[Vista 3.4.6]{Weibel1994HomologicalAlgebra}。
